\documentclass[../../../../SoftwareRequirementsSpecification.tex]{subfiles}
\begin{document}
% Richiede le macro definite in src/macros/useCaseMacros.tex
% (incluse nel preambolo di SoftwareRequirementsSpecification.tex)

\ucstart
\begin{xltabular}{\textwidth}{|l|X|}
  \hline
  \ucid{A-01} \\ \hline
  \ucname{Completamento Profilo} \\ \hline
  \ucactor{Atleta} \\ \hline
  \ucdescription{Al primo accesso l'atleta sceglie una propria
    password, in sostituzione di quella temporanea, e inserisce i
    dati anagrafici e fisici che il PT non ha fornito al momento
    della creazione dell'account.} \\ \hline
  \ucprecond{L'atleta è autenticato (U-01) con la password
    temporanea ricevuta via email. Il profilo dell'atleta è incompleto: il
      PT ha
    creato l'account (PT-02) fornendo solo nome, cognome ed
    email.} \\ \hline
  \ucpostcond{L'atleta ha una propria password e la password
    temporanea non è più valida. Il profilo dell'atleta è completo e
    l'atleta può usare tutte le funzioni dell'applicativo.} \\ \hline
  \ucmainflow{
    \item Il sistema rileva che il profilo dell'atleta è incompleto.
    \item Il sistema mostra il form con i dati mancanti: la nuova
      password e i dati anagrafici e fisici (data di nascita,
      altezza, peso).
    \item L'atleta inserisce i dati.
    \item L'atleta preme il pulsante "Completa Registrazione".
    \item Il sistema valida i dati inseriti.
    \item Il sistema salva i dati nel database, sostituendo
      l'impronta della password temporanea con quella della password
      scelta dall'atleta.
    \item Il sistema indirizza l'atleta alla propria area personale.
  } \\ \hline
  \ucaltflow{
    \textbf{2a. L'atleta abbandona prima di completare:} \newline
    - Il profilo resta incompleto e nessun dato viene salvato. La
    password temporanea resta valida. \newline
    - Al successivo accesso il sistema ripropone il flusso dal
    punto 1. \newline
    \textbf{5a. Dati non validi:} \newline
    - Il sistema mostra un messaggio di errore ("Controllare i dati
    inseriti") e indica i campi non validi (password che non rispetta
    i requisiti minimi di lunghezza e composizione, data di nascita
    futura, altezza o peso non positivi). \newline
    - Il flusso riprende dal passo 2.
  } \\ \hline
  \ucnote{
    Finché il profilo è incompleto, l'atleta non può accedere alle
    altre funzioni dell'applicativo. I dati richiesti sono
    obbligatori: la data di nascita alimenta l'operazione
    \texttt{getAge()} di \texttt{Athlete}, mentre altezza e peso
    servono al PT per stabilire i carichi. L'atleta potrà
    modificare questi dati in seguito, con il caso d'uso A-04
    (Modifica Dati Anagrafici). \newline
    Il cambio della password sta qui e non in un caso d'uso separato
    perché condivide la stessa precondizione: è l'unica cosa che
    l'atleta può fare al primo accesso. La password temporanea
    generata in PT-02 è nota a chiunque abbia letto l'email, quindi
    va sostituita prima di dare all'atleta accesso ai propri dati.
  } \\ \hline
\end{xltabular}
\end{document}
